\section{Divisibility of the low-degree determinants}\label{sec:content}
The full radical in Section~\ref{sec:compression} forces degree $K$.
We now omit it and retain only the added low-degree directions. Their
number is $d=\lfloor K/5\rfloor$. This changes both the determinant and
its arithmetic; savings for the former full matrix cannot simply be
transferred to this one.

\subsection{The fixed lift and the statement}
Use the $J_m$ basis of Section~\ref{sec:moments}. For the same block
zero $z$ from \eqref{eq:Zheight}, define
\begin{equation}\label{eq:low-lift}
 U_{mi}=(-1)^{m-1}
 \frac{(K-a)!}{(K-m)!}\frac{(K+b)!}{(K+m)!}z_m,
 \quad a=5i+1,\ b=a+4,
\end{equation}
and set the other entries of column $i$ to zero. The factorial ratios
contain four factors, each at most $2K$. Thus
\begin{equation}\label{eq:Wlow}
 |U_{mi}|\leq W_K=(2K)^4[15K(4K)^8]^2.
\end{equation}
For $P_i=\sum_mU_{mi}J_m$, direct evaluation gives
\[
 m^2P_i(-4m^2)=\tfrac12\,4^{K-1}(K-a)!(K+b)!z_m
\]
on its block. The same isotropic equation therefore gives $U^TC_KU=0$.
The columns are independent and integral. Set
\[
 G_U(X)=U^T(R_K+XB_K)U,\quad D_U(X)=\det G_U(X),
 \quad F_U=q_UD_U,\quad A_U=q_U|\ell_U|,
\]
where $F_U$ is primitive and $\ell_U$ is the highest nonzero coefficient.
The actual matrix is positive. We do not assume that $\deg D_U=d$.

\begin{theorem}\label{thm:content}
Uniformly over the block choices satisfying
\eqref{eq:Zheight}--\eqref{eq:low-lift},
\begin{equation}\label{eq:arith-rate}
 \limsup_{K\to\infty}\frac{\log A_U}{K^2}
 \leq c_{\rm ar}=\frac{197+192\log2}{150}
 =2.2005617244\ldots.
\end{equation}
The proof gives a finite coefficient-content divisor and an effective
bound before taking this limit.
\end{theorem}
The argument has three distinct content gains. We describe them
separately before combining them, so overlapping prime factors are
not counted twice.

\subsection{A global entry clearer, with the tails retained}
The useful explicit odd residue is
\begin{equation}\label{eq:rhoformula}
 \rho_{2j+1}=
 \frac{2(K-j)(K+j+1)}{4^{2K}}
 \binom{2K-2j}{K-j}\binom{2K+2j+2}{K+j+1}.
\end{equation}
In particular it has no odd-prime denominator. Write
$T_l=\rho_lV_{\cdot l}V_{\cdot l}^T$.

For $p>2K$, collisions of the nodes $1,\ldots,2K$ occur only in
even--odd pairs with sum $p$, if they occur at all. Pairing their
residues gives, over the localized rationals at $p$,
\begin{equation}\label{eq:largep}
 pM_K(X,Z)\equiv
 \left(4\chi X-\frac7{32}Z+\frac14
 \sum_{v=1}^{(p-1)/2}\frac{(-1)^{v-1}}{v^4}\right)N_p\pmod p,
 \qquad \chi=(-1)^{(p-1)/2}.
\end{equation}
To check the common scalar, put $j=(p-1)/2-m$ in an even--odd pair.
The rational part of the resolvent difference, after removing its
common signed cubic factor, is
\[
 \tfrac14\sum_{v=1}^{m}\frac{(-1)^{v-1}}{v^4}
       -4\chi\sum_{v=0}^{j-1}\frac{(-1)^v}{(2v+1)^4}.
\]
Adjacent values agree modulo $p$, since the two new denominators
are opposite up to a factor two; at $m=(p-1)/2$ it is the sum in
\eqref{eq:largep}. The two rational endpoint corrections in the
resolvents are what extend the even sum through $m$.
All these tail denominators are below $p$. At paired roots, the
two values $p/D'$ are negatives modulo $p$, while the numerators
agree. This proves \eqref{eq:largep}. Its $Z$ coefficient is a
$p$-unit. Exact isotropy forces $U^TN_pU=0$ modulo $p$, so every
coefficient of $G_U$ is $p$-integral. This conclusion applies after
cancellation, not to arbitrary raw mixed minors.

For $K<p\leq2K$, set
\[
 Y_p=\sum_{p<l<2K,\ l\ \mathrm{odd}}b_lT_l,
 \qquad R_K=R_{0,p}-\chi p^{-4}Y_p.
\]
The decomposition removes the only $p$-divisible denominator in
each affected odd tail. The matrices $p^2R_{0,p}$, $p^2B_K$,
$p^2Y_p$ and $p^2C_K/7$ are $p$-integral, and
\begin{equation}\label{eq:pwindow}
 p^2Y_p\equiv64\chi\,p^2C_K/7\pmod p.
\end{equation}
Here division by seven is performed over $\Q$ before reduction, so
the formula includes $p=7$.
For completeness, a possible upper odd node is $p+2v$. Its class
contains $2v,p-2v,p+2v$, with a fourth node $2p-2v$ when $v\geq p-K$.
Four-node classes have residue valuation at least $-1$. In a
three-node class only the entry $(v,v)$ can have valuation $-2$.
The two local root-difference products, divided by $p^2$, are
$-16v^2$ at $2v$ and $32v^2$ at $p+2v$. Their residue ratio is
$-1/2$. Moreover $b_{p+2v}\equiv4\chi(-1)^vv^3$ and
$c_v/7=(-1)^{v-1}v^3/32$, giving \eqref{eq:pwindow}.
All other residues disappear in this reduction. Compressing and using
isotropy shows that $pY_{p,U}$, $p^5R_U$ and $p^2B_U$ are integral
at $p$.

For the remaining sufficiently large odd primes, $p^2>4K$, six powers
suffice even before compression:
\begin{equation}\label{eq:p6}
 p^6R_K,\quad p^2B_K,\quad p^2C_K
 \quad\hbox{are $p$-integral}.
\end{equation}
Every central binomial in \eqref{eq:e} has valuation at most one,
and every tail has valuation at least $-4$. If $p\nmid m$, this
gives $v_p(e_m)\geq-2$, $v_p(r_m^e)\geq-4$. If $p\mid m$,
the bounds are $-4$ and $-1$: the cubic numerator pays three tail
powers and the endpoint costs one. For an odd node $l$ with
$p\nmid l$, each gap $(2m-l)(2m+l)$ costs at most one power, giving
$v_p(T_l)\geq-2$. If $p\mid l$, two powers per gap are possible,
but $v_p(b_l)=3$ and $v_p(r_l^o)\geq-1$. These four cases prove
\eqref{eq:p6}, including the endpoint terms.

Put $k_p=\lfloor\log_p(4K)\rfloor$ and
\begin{equation}\label{eq:Qclear}
 Q_K=2^{4K+4+8k_2}
 \prod_{\substack{3\leq p\leq K\\p^2\leq4K}}p^{8k_p}
 \prod_{\substack{3\leq p\leq K\\p^2>4K}}p^6
 \prod_{K<p\leq2K}p^5,
\end{equation}
with products over primes. Then $Q_KR_U,Q_KB_U$ are integral.
For small odd primes, each of the two binomial valuations and each
linear gap costs at most $k_p$; the tails cost at most $4k_p$.
Thus $8k_p$ suffices. At two the odd gaps are units, the denominator
of $\rho_l$ costs at most $4K-1$, and the finite tails and even
terms are covered by $4+8k_2$. This proves the remaining factors.
The prime number theorem gives
\begin{equation}\label{eq:Qrate}
 \log Q_K=(11+4\log2+o(1))K.
\end{equation}
The small primes cost $O(\sqrt K\log K)$; the two large-prime ranges
have weights six and five.

\subsection{Exact tail rank gives three powers per missing direction}
For $K<p\leq2K$, the number of terms in $Y_p$ is
$t_p=K-(p+1)/2$. They are rank one over $\Q$. The exact identity
\[
 p^5G_U(X)=-\chi A_p+p^3S_p(X),\qquad
 A_p=pY_{p,U},\quad S_p=p^2(R_{0,p,U}+XB_U)
\]
has $p$-integral entries and $\rank A_p\leq t_p$ over $\Q$.
In the column expansion of the determinant, more than $t_p$ columns
from $A_p$ give zero identically. Every surviving term therefore
contains at least $(d-t_p)_+$ columns from $p^3S_p$. It follows that
\begin{equation}\label{eq:tailcontent}
 \mathcal T_K=\prod_{K<p\leq2K}p^{3(d-t_p)_+}
 \quad\hbox{divides }\cont(Q_K^dD_U),\qquad
 \log\mathcal T_K=(3/25+o(1))K^2.
\end{equation}
The rate is $3\int_{8/5}^2(x/2-4/5)\dd x=3/25$.
Rank only modulo $p$ would not justify three powers per direction;
the exact rational rank is what makes that conclusion possible.

\subsection{A common dyadic factor}
Set $\eta_0=\eta_1=0$ and
\begin{equation}\label{eq:eta}
 \eta_s=\sum_{j=1}^{s-1}(8j-1-2s_2(j)),\qquad s\geq2.
\end{equation}
Every $s$-minor of $V_{ml}=(4m^2-l^2)^{-1}$ has dyadic valuation
at least $\eta_s$. The Cauchy determinant denominator is odd.
For its numerator use, for $j\geq1$, the integer-valued polynomials
\[
 \frac{\prod_{i=0}^{j-1}(t-i^2)}{j(2j-1)!},\qquad
 \frac{\prod_{i=0}^{j-1}(t-(2i+1)^2)}{4^j(2j)!}.
\]
Take the degree-zero member of each sequence to be $1$.
At $t=m^2$ and $t=(2k+1)^2$ their values are respectively
$\binom{m+j}{2j}+\binom{m+j-1}{2j}$ and $\binom{k+j}{2j}$.
Evaluation determinants therefore divide the even and odd
Vandermondes by the factors
$\prod_{j=1}^{s-1}4^j j(2j-1)!$ and
$\prod_{j=1}^{s-1}4^j(2j)!$.
Their total dyadic valuation is \eqref{eq:eta}.

Cauchy--Binet preserves this bound in $U^TV$. Write the pencil
scaled by $2^{4K+4+8k_2}$ as the sum of its odd-residue Gram part
and its even rational part. The first has integral diagonal weights,
and every $s$-minor is divisible coefficientwise by $2^{2\eta_s}$.
The formulas \eqref{eq:e}, \eqref{eq:evencoeff}, and
$v_2\binom{2n}{n}=s_2(n)$ show that every entry of the even part
has valuation at least $8K-2$. Since $d\leq K/5$,
\[
 2(\eta_d-\eta_s)\leq16(d-1)(d-s)\leq(8K-2)(d-s).
\]
The complementary-minor expansion consequently proves
\begin{equation}\label{eq:dyadiccontent}
 2^{2\eta_d}\mid\cont(Q_K^dD_U),\qquad
 2\eta_d\log2=(8\log2/25+o(1))K^2.
\end{equation}
This applies to every coefficient, not just the leading one.

\subsection{Two valuation layers below the pole cutoff}
Fix an odd prime $p\leq K$ with $p^2>4K$. In a nonzero square
class of nodes $l^2\bmod p$, let $e,o$ count even and odd nodes,
and let $\epsilon_m$ indicate whether $2m$ belongs to the class.
Every root difference within that class has valuation exactly one.
Counting factors in $J_mJ_n/D'$ gives the residue valuations
\begin{equation}\label{eq:classval}
 e-o+1-\epsilon_m-\epsilon_n
 \quad\hbox{at odd nodes},\qquad
 e-o-1\quad\hbox{at an even diagonal node}.
\end{equation}

Write $n=\lfloor2K/p\rfloor$, $s=2K-np$ and $a=\min(s,p-s-1)$.
Represent a nonzero class by $1\leq r<p/2$. Its $2n$ base nodes
$r+jp,p-r+jp$ ($0\leq j<n$) are balanced in parity. Exactly one
extra node occurs precisely for $1\leq r\leq a$; its parity is
that of $n+r$. Let $I_p$ be the even cofactor indices in
odd-majority classes and $Z_p$ those in even-majority classes.
Zero-class indices belong to neither set. If $N_I,N_Z$ count the
two parities of $r$, then
\[
 |I_p|=nN_I,\quad |Z_p|=(n+1)N_Z,\quad
 N_I+N_Z=a,\quad |N_I-N_Z|\leq1.
\]
The sets are unions of at most $n$ and $n+1$ integer intervals.
For example, in $2K/3<p\leq K$ a class first has four balanced
nodes; a fifth makes it unbalanced and a sixth restores balance.
Only odd-majority classes contribute the worst rational layer.

Finite tails cost at most four powers. Formula \eqref{eq:classval}
therefore gives $v_p(R_{mn})\geq-6$ everywhere, at least $-5$ if
one index is outside $I_p$, and at least $-4$ if one index is in
$Z_p$ and the other is outside $I_p$. The zero class needs a separate
check: it consists of $p,2p,\ldots,np$, with squared differences
of valuation two. Double the counts in \eqref{eq:classval}; the
cubic resolvent factor pays three tail powers. All its rational
contributions are at least $-5$, and those involving a $Z_p$ index
are at least $-3$. Thus they respect the stated bounds. The beta
matrix has valuation at least $-2$ and does not weaken them.

Let $b_p$ count fixed five-node blocks meeting $I_p$, let $z_p$ count
blocks wholly in $Z_p$, and put $c_p=d-b_p-z_p$. After compression,
the coefficientwise lower valuation table of $p^6G_U$ is
\[
 \begin{array}{c|ccc}
 &I&\text{remaining}&Z\\\hline
 I&0&1&1\\
 \text{remaining}&1&1&2\\
 Z&1&2&2
 \end{array}.
\]
Divide every $Z$ row and column by $p$, and then every remaining
row by $p$. The result is still integral at $p$. Hence
\begin{equation}\label{eq:Gamma}
 \Gamma_K=\prod_{\substack{p\leq K,\ p\ \mathrm{odd}\\p^2>4K}}
       p^{d-b_p+z_p}\quad\hbox{divides }\cont(Q_K^dD_U).
\end{equation}
This exponent already includes the first-layer rank defect $d-b_p$;
that factor must not be multiplied in again.

Interval counting gives, uniformly,
\[
 d-b_p+z_p=K/5+a/10+O(K/p).
\]
The weighted error sums to $o(K^2)$. For $x=p/K$, the continuous
version of $a/K$ is
$f(x)=\min(2-nx,(n+1)x-2)$ with $n=\lfloor2/x\rfloor$.
Its graph on $2/(n+1)<x\leq2/n$ is a triangle. Therefore
\begin{align}
 \lim K^{-2}\log\Gamma_K
 &=\frac15+\frac1{10}\int_0^1f(x)\dd x\nonumber\\
 &=\frac15+\frac1{10}\sum_{n=2}^{\infty}
                       \frac2{n(n+1)(2n+1)}
 =\frac{23}{30}-\frac45\log2.\label{eq:Gammarate}
\end{align}
Partial summation uses the prime number theorem on
$\varepsilon K<p\leq K$; the omitted content is at most
$2d\vartheta(\varepsilon K)$, so $\varepsilon\downarrow0$ justifies
the endpoint. The series follows by partial fractions and harmonic
number differences.

\subsection{Primitive arithmetic and the compatible comparison}
The three content factors \eqref{eq:tailcontent},
\eqref{eq:dyadiccontent}, and \eqref{eq:Gamma} have disjoint prime
supports. To bound coefficient sizes, use \eqref{eq:rhoformula} and
$\binom{2n}{n}\leq4^n$ to get $\rho_l\leq8K(K+1)\leq16K^2$.
The alternating finite tails have absolute value at most one, so
$|b_l|\leq4K^3$ and $|r_l^o|\leq5K^3$. Since every gap has magnitude
at least one, the total odd contribution is at most $64K^6$ in $B_K$
and $80K^6$ in $R_K$. The reverse binomial bound gives
\[
 |e_m|\leq\frac{(2K+1)^2}{8m^2},\qquad
 |r_m^e|\leq\frac{3m^3}{8},
\]
so the even contribution is at most $K^3$. Therefore
$|R_{mn}|,|B_{mn}|\leq81K^6$, and every
compressed entry has magnitude at most
$M_K=2025W_K^2K^6$, and every determinant coefficient is bounded by
$2^dd!M_K^d$. It follows that
\begin{equation}\label{eq:finalfinite}
 A_U\leq
 \frac{Q_K^d\,2^dd!M_K^d}{\mathcal T_K\,2^{2\eta_d}\Gamma_K}.
\end{equation}
This includes primes from the numerator of $\ell_U$, regardless of
its degree. Combining the rates proves Theorem~\ref{thm:content}.

To interpret it, put $\delta_U=D_U(\bet)/|\ell_U|$. An
irrationality argument would follow from $A_U\delta_U\leq
e^{-\varepsilon K^2}$ along an unbounded sequence. The available
analytic restriction is
\begin{equation}\label{eq:ceiling}
 \limsup\frac{-\log\delta_U}{K^2}
 \leq\frac{\pi/2+2\log(5/4)+2\arctan2+4\log(2+\sqrt5)}5
 =2.0011845531\ldots.
\end{equation}
Appendix~\ref{app:ceiling} proves it for every degree-below-$K$
subspace of dimension at most $K/5$, without a beta-rank assumption.
The arithmetic upper bound $2.20056\ldots$ is above this ceiling.
Thus these particular estimates cannot establish the desired
comparison. They do not supply a necessary arithmetic cost, so they
do not prove that the primitive determinants themselves fail.
